\chapter{Zaključak}
\label{ch:zakljucak}

U radu je razvijeno i primijenjeno okruženje za mjerenje promjene govornog signala pri transkodiranju u FreeSWITCH-u. Glavna vrijednost pristupa jeste izdvajanje ulaznog kanala A i izlaznog kanala B neposredno iz RTP paketa, čime se posmatra obrada na medijskom serveru bez oslanjanja na njegove interne PCM snimke. Dopunsko poređenje originalnog WAV-a i B-lega pokazuje konačnu očuvanost signala, a koristi iste snimljene pozive bez pokretanja novih testova.

\section{Odgovori na istraživačka pitanja}

Za punu matricu pet kodeka prosječni PESQ MOS-LQO iznosi \PassthroughPESQ{} bez transkodiranja i \TranscodePESQ{} s transkodiranjem. Kada se svaka konfiguracija usporedi sa kontrolom istog izvornog kodeka, prosječna dodatna degradacija iznosi \RazlikaPESQ{} boda. Nalaz je potvrđen eksplorativnom analizom konfiguracijskih razlika, ali ne znači da svaki smjer gubi isti iznos.

Smjer i odredišni kodek imaju velik utjecaj. Kodiranje u Opus uglavnom dobro čuva signal kanala A, dok kodiranje u GSM dodaje najveću degradaciju. Usporedba svih deset parova pokazuje da suprotni smjerovi nisu zamjenjivi, zbog različitog odredišnog kodera i promjene frekvencije uzorkovanja.

Prosječno opterećenje kontejnera približno je \OmjerCPU{} puta veće u grupi s transkodiranjem. Najviše vrijednosti zabilježene su pri kodiranju u Opus. Taj nalaz predstavlja relativnu razliku na ispitnom računaru i nije dovoljan za neposredno dimenzioniranje produkcijskog sistema.

End-to-end analiza potvrđuje da se mala dodatna degradacija ne smije izjednačiti s visokim konačnim kvalitetom. Prosječni WB PESQ iznosi 3,170 za kontrole i 2,608 za transkodirane konfiguracije, uz prosječnu dodatnu degradaciju od 0,563 boda prema kontroli istog izvora. Putanja GSM$\rightarrow$PCMA to jasno pokazuje: A-leg$\rightarrow$B-leg rezultat je 4,172, a rezultat prema originalnom WAV-u 2,255. FreeSWITCH u toj putanji dodaje samo 0,039 boda degradacije, ali ne može nadoknaditi informaciju koju je početni GSM koder već odbacio. Spektralna analiza smjera Opus$\rightarrow$GSM dodatno potvrđuje gubitak viših frekvencija: udio energije u opsegu 3,4--7 kHz smanjen je sa 2,84\% na 0,08\%.

\section{Ograničenja zaključaka}

Mjerenje između kanala A i kanala B izoluje doprinos servera, dok mjerenje od originalnog WAV-a obuhvata cijeli posmatrani kodečki lanac. Nijedno samo po sebi nije potpuna procjena korisničkog iskustva u glasovnom pozivu. Zajednički WB postupak na 16 kHz može utjecati na uskopojasne konfiguracije, pa dopunska analiza prikazuje i PESQ u izvornom NB ili WB režimu. Posebno ponašanje konfiguracija s Opusom pokazuje da na rezultat mogu utjecati rekonstrukcija RTP toka, Ogg wrapper i vremensko poravnanje.

Korišten je jedan sintetizirani glas, bez namjerno uvedenog gubitka paketa i kašnjenja. Eksperiment je proveden na jednom računaru, a CPU mjerenje obuhvata cijeli kontejner, osnovno opterećenje i dio pripremnog intervala. Deset ponavljanja po konfiguraciji procjenjuje tehničku ponovljivost, ali ne zamjenjuje raznolik skup govornika ni subjektivni MOS test prema ITU-T P.800.

\section{Preporuke i budući rad}

U upravljanim VoIP mrežama treba dati prednost zajedničkom kodeku na oba kanala. Kada je transkodiranje potrebno, politike treba zasnivati na konkretnom smjeru, a ne samo na nazivima dva kodeka. Za promet prema Opusu treba planirati veće procesorske resurse, dok GSM kao odredišni kodek treba koristiti kada kompatibilnost ili ušteda propusnosti opravdavaju izmjerenu degradaciju.

Buduća provjera treba koristiti više stvarnih govornika oba spola i ujednačen skup rečenica na jeziku ciljnih korisnika. Potrebno je dodatno provjeriti rekonstrukciju Opus RTP toka. CPU treba mjeriti samo tokom stabilnog aktivnog dijela poziva i nakon oduzimanja osnovnog opterećenja. Ako su dostupne odgovarajuća licenca i referentna implementacija, rezultate treba potvrditi metrikom POLQA. Zaseban eksperiment može kontrolisano uvesti gubitak paketa i kolebanje kašnjenja, porediti različite mehanizme prikrivanja gubitka te uključiti subjektivno slušno ispitivanje.

Zaključno, implementirani sistem daje ponovljiv način poređenja svih usmjerenih kombinacija odabranih kodeka na nivou poslužitelja i cijelog kodečkog lanca. Najvažniji nalaz nije jedna univerzalna vrijednost degradacije, nego činjenica da učinak transkodiranja zavisi od smjera, odredišnog kodeka, prethodno izgubljenog sadržaja, procesorskog troška i izbora referentne tačke.
